% CLIMB — figure captions. Generated 2026-08-20.
% Six canonical panels throughout: MoleculeACE (30 ChEMBL targets, macro RMSE), HIV (NEF1\%),
% BACE, Ames and Tox21 (ROC-AUC), QM7 (RMSE). Arrows in panel titles give the good direction.

% ---------------------------------------------------------------- Fig A
\caption{\textbf{Overall standing across 65 datasets.}
\textbf{(a)} Each model is ranked within every dataset and the ranks are averaged per suite;
the plotted value is the mean of the four suite means, so MoleculeNet (6), MoleculeACE (30),
Polaris (28) and CBS (1) each count once rather than letting the two large suites decide the
order. Filled dot: mean rank ($\pm$SE across suites); open markers: the four suite means;
$\times$: the model is fine-tuned end-to-end rather than probed frozen. Arms are admitted at
$\ge$60 of 65 datasets. Each representation is shown at the probe head that suits it
(Supp.~Fig.~f).
\textbf{(b)} Absolute performance on the six canonical panels. Error bars are 95\% cluster
bootstrap intervals over test scaffolds (targets on MoleculeACE); Ames is analytic
(Hanley--McNeil), since its test labels are withheld. Dotted line: untrained (random-init,
frozen) encoder.
Classical fingerprint baselines lead overall, with the CheMeleon foundation model next and the
CLIMB encoders behind both.}

% ---------------------------------------------------------------- Fig B
\caption{\textbf{Pretraining scaling.}
Downstream performance against pretraining tokens actually processed (log scale), one line per
pretraining recipe, seed-0 rungs. Reference lines: the two XGBoost fingerprint baselines
(dash-dot, dashed) and the untrained encoder (dotted). The top two unsupervised rungs (open
markers) use a larger RDKit-canonical corpus than the rungs below them.
Scaling the pretraining corpus by two orders of magnitude does not carry any recipe past the
fingerprint baselines on these tasks.}

% ---------------------------------------------------------------- Fig C+D
\caption{\textbf{Does similarity to the pretraining data explain the benefit?}
Lift is measured against an untrained, frozen encoder probed identically, so $0$ means the
pretraining bought nothing.
\textbf{(a--c)} Unsupervised pretraining: per-molecule lift against maximum ECFP4 Tanimoto
similarity to the pretraining corpus. Gains are flat in similarity --- molecules chemically
close to the corpus benefit no more than distant ones.
\textbf{(d--f)} Supervised pretraining: each point is one (SFT label family $\times$ evaluation
task) pair; $x$ is mean max Tanimoto similarity between the family's molecules and the task's
(a lower bound: families are sampled). Transfer is likewise unexplained by chemical overlap,
which points to the label type rather than the molecules as what carries it.}

% ---------------------------------------------------------------- Fig E
\caption{\textbf{Corrupted pretraining objectives.}
Objective family, data volume, compute, schedule, architecture and probe are held fixed; only
the chemical content of the pretraining signal is destroyed. Bars are lift over an untrained,
frozen encoder.
\textbf{(a)} Supervised descriptor regression with targets permuted across the batch: the
molecule$\rightarrow$descriptor mapping is destroyed, the target distribution is not.
\textbf{(b)} The unsupervised objective with its chemical content removed.
The real objectives help on every panel and the corrupted controls fall below the untrained
floor, so the gains do come from chemistry rather than from optimisation alone --- but they
start from, and return to, a low baseline.}

% ---------------------------------------------------------------- Fig F
\caption{\textbf{Are the learned embeddings redundant given classical features?}
The same XGBoost, the same scaffold 5-fold CV and the same task, given three classical feature
bases (descriptors, ECFP4, both) with and without a 512-d frozen embedding concatenated on.
Raw features are concatenated without standardisation (XGBoost is scale-invariant). Error bars
are $\pm1$ SD over the panel's replicate unit --- targets on MoleculeACE, folds elsewhere ---
and analytic (Hanley--McNeil) on Ames, whose test labels are withheld.
Adding either CLIMB embedding to descriptors+ECFP4 fails to beat its own spread on any of the
six panels and costs 0.036 macro RMSE on MoleculeACE and 0.021 ROC-AUC on Ames; adding
CheMeleon on the identical bases gains 0.013 on Ames. The CLIMB representations are therefore
informationally redundant given fingerprints and descriptors, which is a stronger statement
than scoring lower.}

% ---------------------------------------------------------------- Fig G
\caption{\textbf{Does the representation respond when the chemistry changes?}
Ten chemical edits, 100 molecule pairs each, in which the two molecules are genuinely different
and a usable representation must respond. Response is the embedding shift the edit produces
divided by \emph{that model's own} shift when a completely different compound of matched
molecular weight is substituted, so $1.00$ means ``moves the representation as far as changing
the molecule''. Panel (i) is that reference and is $1.000$ by construction. The per-model
denominator lets a 512-d transformer and a 2048-bit fingerprint share one axis. All models
receive RDKit-canonical SMILES.
Notation-only edits (SMILES enumeration, Kekulisation) should score $0$; stereochemical and
scaffold edits should not.}

% ---------------------------------------------------------------- SI Fig a
\caption{\textbf{Is end-to-end fine-tuning necessary?}
The same pretrained encoder used two ways at full downstream data: frozen with a trained probe,
versus the whole network fine-tuned. All points come from one wave, so frozen and end-to-end are
like-for-like within and across encoders in a panel. Intervals are 95\% resampling intervals,
one method per panel: cluster bootstrap over test scaffolds, over targets on MoleculeACE, and
analytic on Ames. Dotted line: XGBoost on ECFP4+descriptors.
End-to-end is ahead in 12 of 18 cells, but no frozen/end-to-end interval pair is disjoint --- at
this test-set size the data does not resolve any single one of these differences. CheMeleon is
shown at its XGBoost probe, where freezing is ahead on four of six panels.}

% ---------------------------------------------------------------- SI Fig b
\caption{\textbf{Tokenizer family and vocabulary size.}
Two tokenizer families (byte-level BPE, Unigram-LM), each at four distinct vocabulary sizes
spanning the reachable range; one MLM encoder per tokenizer at matched forward passes, same
corpus, same frozen probe. $x$ is the vocabulary actually reached (log scale).
Neither the family nor the size moves downstream performance beyond replicate noise.}

% ---------------------------------------------------------------- SI Fig c
\caption{\textbf{Featurization cost.}
Wall-clock to featurize 1000 molecules, 5 repeats after warm-up, with the settings used in
production. The transformer encoder is orders of magnitude more expensive than the classical
featurizers it does not outperform, which is the relevant comparison at virtual-screening
scale.}

% ---------------------------------------------------------------- SI Fig d
\caption{\textbf{Canonical versus enumerated SMILES in pretraining.}
Randomised SMILES are standard practice, intended to teach the model that the string is
arbitrary and the graph is not. Canonical and enumerated pretraining are matched at five corpus
fractions $\times$ 3 pretraining seeds, so augmentation is tested at every rung --- including
the scarce-data rungs where it should help most.
It does not help at any rung.}

% ---------------------------------------------------------------- SI Fig e
\caption{\textbf{Label efficiency.}
Absolute performance against the number of labelled training molecules (log scale), with an
identical hold-out split, label fractions and seed grid across models. This is where a frozen
pretrained encoder is most likely to pay: the small-data regime.}

% ---------------------------------------------------------------- SI Fig f
\caption{\textbf{The probe head is not neutral.}
The same frozen representations scored through an MLP head and through XGBoost, everything else
fixed. The preference is representation-dependent and consistent across panels: the CLIMB
encoders do better with the MLP head, while fingerprints, descriptors and CheMeleon do better
with XGBoost. Comparing representations across different heads therefore confounds the two, and
each representation is reported at its own best head in Fig.~1a.}
